\documentclass{article}
\usepackage{amsmath,amssymb,amsthm}
\usepackage{algorithm}
\usepackage[noend]{algpseudocode}
\usepackage{enumitem}
\usepackage{hyperref}
\usepackage{geometry}
\geometry{margin=1in}
\newtheorem{lemma}{Lemma}
\newtheorem{theorem}{Theorem}
\newtheorem{assumption}{Assumption}
\newcommand{\norm}[1]{\left\lVert#1\right\rVert}
\newcommand{\R}{\mathbb{R}}

\begin{document}

\section*{Algorithm Name}
\textbf{Trust‑Region Method with Approximate Subproblem Solution (TR‑AS)}  

The algorithm iteratively builds a quadratic model of a smooth non‑convex objective
\(f:\R^{n}\rightarrow\R\) around the current iterate, solves a trust‑region subproblem
approximately, and updates both the iterate and the trust‑region radius according to
the quality of the obtained step.

\section*{Mathematical Formulation}
At iteration \(k\) let \(x_k\in\R^{n}\) be the current point and \(\Delta_k>0\) the trust‑region radius.  
Define

\[
g_k \;:=\; \nabla f(x_k), \qquad
B_k \;:=\; \nabla^{2} f(x_k) \;\; \text{(or a symmetric approximation)} .
\]

The quadratic model is

\[
m_k(p) \;=\; f(x_k) + g_k^{\top}p + \tfrac12 p^{\top}B_k p ,\qquad p\in\R^{n}.
\]

The (exact) trust‑region subproblem is

\[
\min_{p\in\R^{n}}\; m_k(p) \quad \text{s.t.}\quad \norm{p}\le \Delta_k .
\tag{TR}
\]

TR‑AS computes an \emph{approximate} solution \(p_k\) that satisfies the
\textbf{Cauchy decrease condition}

\[
m_k(0)-m_k(p_k) \;\ge\; \kappa_{\mathrm{c}}\;
\bigl(m_k(0)-m_k(p_k^{\mathrm{C}})\bigr),
\tag{1}
\]
where \(p_k^{\mathrm{C}}\) is the Cauchy point (the minimizer of \(m_k\) along the
steepest‑descent direction) and \(\kappa_{\mathrm{c}}\in(0,1]\) is a prescribed constant.

Define the \emph{actual reduction} and \emph{predicted reduction}

\[
\operatorname{ared}_k \;:=\; f(x_k)-f(x_k+p_k),\qquad
\operatorname{pred}_k \;:=\; m_k(0)-m_k(p_k).
\]

The ratio

\[
\rho_k \;:=\; \frac{\operatorname{ared}_k}{\operatorname{pred}_k}
\tag{2}
\]

measures the quality of the step.  Using user‑specified parameters
\(0<\eta_1<\eta_2<1\) (typically \(\eta_1=0.1,\;\eta_2=0.75\)),
the update rules are

\[
x_{k+1}= 
\begin{cases}
x_k+p_k, & \rho_k\ge \eta_1,\\[2mm]
x_k, & \rho_k<\eta_1,
\end{cases}
\tag{3}
\]

\[
\Delta_{k+1}= 
\begin{cases}
\displaystyle \min\{\gamma_{\mathrm{inc}}\Delta_k,\;\Delta_{\max}\},
& \rho_k\ge \eta_2,\\[2mm]
\displaystyle \gamma_{\mathrm{dec}}\Delta_k,
& \rho_k<\eta_1,\\[2mm]
\Delta_k, & \text{otherwise},
\end{cases}
\tag{4}
\]
with expansion factor \(\gamma_{\mathrm{inc}}>1\), contraction factor
\(0<\gamma_{\mathrm{dec}}<1\), and a maximum radius \(\Delta_{\max}>0\).

\section*{Pseudocode}
\begin{algorithm}[H]
\caption{Trust‑Region Method with Approximate Subproblem Solution (TR‑AS)}
\begin{algorithmic}[1]
\State \textbf{Input:} initial point \(x_0\), initial radius \(\Delta_0>0\),
parameters \(\eta_1,\eta_2,\gamma_{\mathrm{inc}},\gamma_{\mathrm{dec}},\Delta_{\max},\kappa_{\mathrm{c}}\).
\State Set \(k\gets0\).
\While{stopping criterion not satisfied}
    \State Compute \(g_k=\nabla f(x_k)\) and a symmetric matrix \(B_k\) (e.g. Hessian).
    \State Approximately solve \((\text{TR})\) obtaining \(p_k\) that satisfies (1).
    \State Compute \(\operatorname{ared}_k = f(x_k)-f(x_k+p_k)\) and
          \(\operatorname{pred}_k = m_k(0)-m_k(p_k)\).
    \State \(\rho_k \gets \operatorname{ared}_k/\operatorname{pred}_k\).
    \If{\(\rho_k \ge \eta_1\)} 
        \State \(x_{k+1}\gets x_k+p_k\) \Comment{Accept step}
    \Else
        \State \(x_{k+1}\gets x_k\) \Comment{Reject step}
    \EndIf
    \If{\(\rho_k \ge \eta_2\)}
        \State \(\Delta_{k+1}\gets \min\{\gamma_{\mathrm{inc}}\Delta_k,\Delta_{\max}\}\)
    \ElsIf{\(\rho_k < \eta_1\)}
        \State \(\Delta_{k+1}\gets \gamma_{\mathrm{dec}}\Delta_k\)
    \Else
        \State \(\Delta_{k+1}\gets \Delta_k\)
    \EndIf
    \State \(k\gets k+1\)
\EndWhile
\State \textbf{Output:} last iterate \(x_k\).
\end{algorithmic}
\end{algorithm}

\section*{Dry Run Example}
Consider the scalar problem
\[
f(w)=(w-3)^2,\qquad w\in\R.
\]
We use the following hyper‑parameters:
\[
\Delta_0=1,\quad \eta_1=0.1,\quad \eta_2=0.75,\quad
\gamma_{\mathrm{inc}}=2,\quad \gamma_{\mathrm{dec}}=0.5,\quad
\kappa_{\mathrm{c}}=1.
\]
Since \(f\) is quadratic, we set \(B_k = \nabla^2 f(w_k)=2\) for all \(k\).
The Cauchy point for a scalar problem coincides with the exact solution of (TR)
provided \(\norm{p}\le\Delta_k\); therefore the approximate solution will be exact.

\begin{center}
\begin{tabular}{c|c|c|c|c}
\(k\) & \(w_k\) & \(g_k=2(w_k-3)\) & \(p_k\) (solution of (TR)) & \(f(w_{k+1})\)\\\hline
0 & 0 & \(-6\) & 
\(\displaystyle p_0 = \min\!\Bigl\{\frac{-g_0}{B_0},\;\Delta_0\Bigr\}
 = \min\{3,1\}=1\) & 
\(f(1) = (1-3)^2 = 4\)\\[2mm]
1 & 1 & \(-4\) & 
\(p_1 = \min\!\Bigl\{\frac{-g_1}{B_1},\;\Delta_1\Bigr\}
 = \min\{2,2\}=2\) (here \(\Delta_1=\gamma_{\mathrm{inc}}\Delta_0=2\))
 & 
\(f(3)=0\)\\[2mm]
2 & 3 & \(0\) & 
\(p_2 = 0\) (gradient zero) &
\(f(3)=0\)\\
\end{tabular}
\end{center}

\noindent
Step details (first iteration):
\[
\begin{aligned}
\operatorname{pred}_0 &= m_0(0)-m_0(p_0)
 = f(0)-\bigl[f(0)+g_0 p_0+\tfrac12 B_0 p_0^2\bigr] \\
 &= -\bigl[-6\cdot 1+ \tfrac12\cdot2\cdot1^2\bigr]
 = 5 .
\end{aligned}
\]
\[
\operatorname{ared}_0 = f(0)-f(1)=9-4=5,\qquad 
\rho_0=\frac{5}{5}=1\ge\eta_2,
\]
so the step is accepted and the radius is expanded to \(\Delta_1=2\).  
The same computation for the second iteration yields \(\rho_1=1\) and the
algorithm terminates with the global minimiser \(w=3\).

\section*{Assumptions}
\begin{assumption}[Smoothness]\label{ass:smooth}
\(f:\R^{n}\rightarrow\R\) is twice continuously differentiable and its gradient
is Lipschitz continuous with constant \(L>0\):
\[
\norm{\nabla f(x)-\nabla f(y)}\le L\norm{x-y},\qquad\forall x,y\in\R^{n}.
\]
\end{assumption}
\begin{assumption}[Bounded Hessian Approximation]\label{ass:hessian}
For all iterates \(k\) the symmetric matrix \(B_k\) satisfies
\[
\|B_k\|\le M,\qquad \text{for some } M>0.
\]
\end{assumption}
\begin{assumption}[Cauchy Decrease]\label{ass:cauchy}
The approximate step \(p_k\) fulfills (1) with a fixed constant
\(\kappa_{\mathrm{c}}\in(0,1]\).
\end{assumption}
\begin{assumption}[Parameter Bounds]\label{ass:params}
The algorithmic parameters satisfy
\[
0<\eta_1<\eta_2<1,\qquad
\gamma_{\mathrm{inc}}>1,\qquad
0<\gamma_{\mathrm{dec}}<1,\qquad
\Delta_{\max}>0.
\]
\end{assumption}

\section*{Main Convergence Theorem}
\begin{theorem}[Global Convergence and Complexity]\label{thm:main}
Under Assumptions \ref{ass:smooth}–\ref{ass:params}, the sequence
\(\{x_k\}\) generated by TR‑AS satisfies
\[
\liminf_{k\to\infty}\norm{\nabla f(x_k)} = 0.
\]
Moreover, for any tolerance \(\varepsilon>0\) the algorithm produces an iterate
\(x_{k}\) with \(\norm{\nabla f(x_k)}\le\varepsilon\) after at most
\[
\mathcal{O}\!\Bigl(\frac{M^2 (f(x_0)-f_{\inf})}{\kappa_{\mathrm{c}}^2 \eta_1^2\,
\varepsilon^{2}}\Bigr)
\]
successful iterations, where \(f_{\inf}:=\inf_{x}f(x)\).
\end{theorem}

\section*{Theoretical Properties}
\begin{itemize}[leftmargin=*]
\item \textbf{Stability of the radius.}  The radius \(\Delta_k\) remains bounded
away from zero as long as \(\norm{\nabla f(x_k)}\) is bounded away from zero;
consequently the algorithm cannot stall prematurely.
\item \textbf{Hyper‑parameter sensitivity.}
Increasing \(\gamma_{\mathrm{inc}}\) or decreasing \(\gamma_{\mathrm{dec}}\) leads to
more aggressive radius adaptation, which may reduce the number of
unsuccessful steps but can increase the risk of overshooting in highly
non‑convex regions.  The constants \(\eta_1,\eta_2\) control the acceptance
threshold; standard choices (\(\eta_1=0.1,\eta_2=0.75\)) guarantee the bounds
used in the proof.
\item \textbf{Comparison to baseline methods.}
Compared with (unconstrained) gradient descent, TR‑AS enjoys a guaranteed
decrease even when the gradient points toward a region of high curvature,
because the trust‑region constraint prevents overly large steps.
The iteration complexity \(\mathcal{O}(\varepsilon^{-2})\) matches that of
gradient descent for smooth non‑convex problems, while often requiring fewer
function evaluations in practice due to better step acceptance.
\end{itemize}

\section*{Discussion}
The theorem shows that the algorithm converges to first‑order stationary points
without assuming convexity.  The complexity bound is derived from a uniform
lower bound on the predicted reduction (Lemma~\ref{lem:pred-lower}) and a
uniform upper bound on the actual reduction error (Lemma~\ref{lem:ared-pred}).
Both lemmas rely only on the Lipschitz continuity of the gradient and the
Cauchy‑decrease condition (1).  Consequently, the result holds for any
approximate subproblem solver that respects (1), including truncated
conjugate‑gradient or dogleg methods.

\section*{Preliminaries (Definitions and Lemmas)}
\begin{lemma}[Cauchy point reduction]\label{lem:cauchy}
Let \(p_k^{\mathrm{C}} = -\frac{g_k^{\top}g_k}{g_k^{\top}B_k g_k}\,g_k\) be the
Cauchy point (projected onto the ball of radius \(\Delta_k\)).  Then
\[
m_k(0)-m_k(p_k^{\mathrm{C}})
\;\ge\;
\frac{1}{2}\,\frac{\|g_k\|^2}{\|B_k\|}\,
\min\!\Bigl\{1,\;\frac{\|g_k\|}{\|B_k\|\Delta_k}\Bigr\}.
\tag{5}
\]
\end{lemma}
\begin{proof}
The quadratic model along the direction \(-g_k\) is
\[
\phi(\alpha)=f(x_k)-\alpha \|g_k\|^2 + \tfrac12\alpha^2 g_k^{\top}B_k g_k .
\]
Minimising \(\phi\) w.r.t. \(\alpha\) yields
\(\alpha^{\star}= \|g_k\|^2/(g_k^{\top}B_k g_k)\).
If \(\alpha^{\star}\Delta_k\le\Delta_k\) the Cauchy point lies inside the ball,
otherwise it is truncated to \(\Delta_k\).  Plugging the two cases into
\(\phi(0)-\phi(\alpha)\) gives (5).  \(\square\)
\end{proof}

\begin{lemma}[Predicted reduction lower bound]\label{lem:pred-lower}
If \(p_k\) satisfies the Cauchy decrease condition (1), then
\[
\operatorname{pred}_k
= m_k(0)-m_k(p_k)
\;\ge\;
\kappa_{\mathrm{c}}\,
\frac{1}{2}\,\frac{\|g_k\|^2}{M}\,
\min\!\Bigl\{1,\;\frac{\|g_k\|}{M\Delta_k}\Bigr\}.
\tag{6}
\]
\end{lemma}
\begin{proof}
From Lemma~\ref{lem:cauchy} and the bound \(\|B_k\|\le M\) (Assumption~\ref{ass:hessian}),
\[
m_k(0)-m_k(p_k^{\mathrm{C}})
\;\ge\;
\frac12\frac{\|g_k\|^2}{M}
\min\!\Bigl\{1,\frac{\|g_k\|}{M\Delta_k}\Bigr\}.
\]
Multiplying by \(\kappa_{\mathrm{c}}\) and using (1) yields (6). \(\square\)
\end{proof}

\begin{lemma}[Actual vs. predicted reduction]\label{lem:ared-pred}
Under Lipschitz gradient (Assumption~\ref{ass:smooth}), the model error satisfies
\[
|\,\operatorname{ared}_k-\operatorname{pred}_k\,|
\;\le\;
\frac{L}{2}\,\norm{p_k}^{3}.
\tag{7}
\]
\end{lemma}
\begin{proof}
By Taylor's theorem with remainder,
\[
f(x_k+p_k) = f(x_k) + g_k^{\top}p_k + \tfrac12 p_k^{\top}B_k p_k
            + R_k,
\]
where \(|R_k|\le \frac{L}{6}\norm{p_k}^{3}\).  Hence
\[
\operatorname{ared}_k
= f(x_k)-f(x_k+p_k)
= -(g_k^{\top}p_k + \tfrac12 p_k^{\top}B_k p_k) - R_k.
\]
Since \(\operatorname{pred}_k = -(g_k^{\top}p_k + \tfrac12 p_k^{\top}B_k p_k)\),
the difference equals \(|R_k|\le \frac{L}{6}\norm{p_k}^{3}\).
A slightly looser bound \(\frac{L}{2}\norm{p_k}^{3}\) is convenient and holds.
\(\square\)
\end{proof}

\section*{Main Proof (Step‑by‑Step Derivation)}
We prove Theorem~\ref{thm:main}.  The proof proceeds by establishing a
uniform lower bound on the predicted reduction, bounding the error between
actual and predicted reduction, and then showing that the trust‑region radius
cannot shrink to zero while the gradient norm stays bounded away from zero.

\noindent\textbf{[Step 1]} (\emph{Bounding the ratio \(\rho_k\)}).  
From Lemma~\ref{lem:ared-pred} we have
\[
\operatorname{ared}_k
\;\ge\;
\operatorname{pred}_k - \frac{L}{2}\norm{p_k}^{3}.
\tag{8}
\]
Dividing (8) by \(\operatorname{pred}_k>0\) yields
\[
\rho_k
\;\ge\;
1 - \frac{L}{2}\,\frac{\norm{p_k}^{3}}{\operatorname{pred}_k}.
\tag{9}
\]

\noindent\textbf{[Step 2]} (\emph{Using the predicted reduction lower bound}).  
Insert (6) into (9).  Since \(\norm{p_k}\le\Delta_k\),
\[
\frac{\norm{p_k}^{3}}{\operatorname{pred}_k}
\;\le\;
\frac{\Delta_k^{3}}
{\kappa_{\mathrm{c}}\frac12\frac{\|g_k\|^{2}}{M}
\min\!\Bigl\{1,\frac{\|g_k\|}{M\Delta_k}\Bigr\}}
=
\frac{2M\,\Delta_k^{3}}
{\kappa_{\mathrm{c}}\|g_k\|^{2}\,
\min\!\Bigl\{1,\frac{\|g_k\|}{M\Delta_k}\Bigr\}}.
\tag{10}
\]

\noindent\textbf{[Step 3]} (\emph{Two cases for the minimum}).  

\begin{enumerate}
\item \textbf{Case A:}\(\displaystyle \frac{\|g_k\|}{M\Delta_k}\ge 1\) (i.e. \(\|g_k\|\ge M\Delta_k\)).  
Then \(\min\{1,\frac{\|g_k\|}{M\Delta_k}\}=1\) and (10) becomes
\[
\frac{\norm{p_k}^{3}}{\operatorname{pred}_k}
\;\le\;
\frac{2M\,\Delta_k^{3}}{\kappa_{\mathrm{c}}\|g_k\|^{2}}
\;\le\;
\frac{2M\,\Delta_k^{3}}{\kappa_{\mathrm{c}}(M\Delta_k)^{2}}
=
\frac{2\Delta_k}{\kappa_{\mathrm{c}}M}.
\tag{11}
\]

\item \textbf{Case B:}\(\displaystyle \frac{\|g_k\|}{M\Delta_k}< 1\) (i.e. \(\|g_k\|< M\Delta_k\)).  
Now the minimum equals \(\frac{\|g_k\|}{M\Delta_k}\) and (10) gives
\[
\frac{\norm{p_k}^{3}}{\operatorname{pred}_k}
\;\le\;
\frac{2M\,\Delta_k^{3}}
{\kappa_{\mathrm{c}}\|g_k\|^{2}\,\frac{\|g_k\|}{M\Delta_k}}
=
\frac{2M^{2}\,\Delta_k^{4}}{\kappa_{\mathrm{c}}\|g_k\|^{3}}.
\tag{12}
\]
Since \(\|g_k\|<M\Delta_k\), we have \(\|g_k\|^{3}<M^{3}\Delta_k^{3}\), thus
\[
\frac{\norm{p_k}^{3}}{\operatorname{pred}_k}
\;<\;
\frac{2M^{2}\Delta_k^{4}}{\kappa_{\mathrm{c}}M^{3}\Delta_k^{3}}
=
\frac{2\Delta_k}{\kappa_{\mathrm{c}}M}.
\tag{13}
\]
\end{enumerate}

\noindent\textbf{[Step 4]} (\emph{Uniform bound on \(\rho_k\)}).  
Both cases yield the same upper bound (11)–(13):
\[
\frac{\norm{p_k}^{3}}{\operatorname{pred}_k}
\;\le\;
\frac{2\Delta_k}{\kappa_{\mathrm{c}}M}.
\]
Plugging into (9) gives
\[
\rho_k \;\ge\;
1 - \frac{L}{2}\cdot\frac{2\Delta_k}{\kappa_{\mathrm{c}}M}
=
1 - \frac{L\Delta_k}{\kappa_{\mathrm{c}}M}.
\tag{14}
\]

\noindent\textbf{[Step 5]} (\emph{Choice of a safe radius}).  
Define
\[
\Delta_{\mathrm{safe}}:=\frac{\kappa_{\mathrm{c}}M}{2L}.
\]
If \(\Delta_k\le \Delta_{\mathrm{safe}}\), then from (14)
\[
\rho_k \ge 1-\frac{L\Delta_{\mathrm{safe}}}{\kappa_{\mathrm{c}}M}
= 1-\frac12 = \frac12 >\eta_1 .
\tag{15}
\]
Hence any iteration with \(\Delta_k\le\Delta_{\mathrm{safe}}\) is \emph{accepted}
and the radius will not be contracted (rule (4)).  Consequently,
\[
\Delta_k \ge \min\{\Delta_{\mathrm{safe}},\Delta_0\}
\quad\text{for all sufficiently large }k.
\tag{16}
\]

\noindent\textbf{[Step 6]} (\emph{Descent when the step is accepted}).  
If \(\rho_k\ge\eta_1\), the actual reduction satisfies
\[
\operatorname{ared}_k
= \rho_k\,\operatorname{pred}_k
\;\ge\;
\eta_1\,\operatorname{pred}_k.
\tag{17}
\]
Using the lower bound (6) we obtain
\[
\operatorname{ared}_k
\;\ge\;
\eta_1\kappa_{\mathrm{c}}\,
\frac12\frac{\|g_k\|^{2}}{M}
\min\!\Bigl\{1,\frac{\|g_k\|}{M\Delta_k}\Bigr\}.
\tag{18}
\]

\noindent\textbf{[Step 7]} (\emph{Summation of reductions}).  
Let \(K\) be the total number of \emph{successful} iterations (those with
\(\rho_k\ge\eta_1\)).  Summing (18) over all successful steps yields
\[
\sum_{k\in\mathcal{S}} \operatorname{ared}_k
\;\ge\;
\frac{\eta_1\kappa_{\mathrm{c}}}{2M}
\sum_{k\in\mathcal{S}} \|g_k\|^{2}
\min\!\Bigl\{1,\frac{\|g_k\|}{M\Delta_k}\Bigr\}.
\tag{19}
\]
Because the objective value is bounded below by \(f_{\inf}\),
\[
\sum_{k\in\mathcal{S}} \operatorname{ared}_k
= f(x_0)-f(x_{K}) \;\le\; f(x_0)-f_{\inf}.
\tag{20}
\]

\noindent\textbf{[Step 8]} (\emph{Deriving a complexity bound}).  
From (19)–(20) we have
\[
\frac{\eta_1\kappa_{\mathrm{c}}}{2M}
\sum_{k\in\mathcal{S}} \|g_k\|^{2}
\min\!\Bigl\{1,\frac{\|g_k\|}{M\Delta_k}\Bigr\}
\;\le\;
f(x_0)-f_{\inf}.
\tag{21}
\]
Using the lower bound on the radius (16), i.e. \(\Delta_k\ge\Delta_{\min}
:=\min\{\Delta_{\mathrm{safe}},\Delta_0\}>0\), we obtain
\[
\min\!\Bigl\{1,\frac{\|g_k\|}{M\Delta_k}\Bigr\}
\ge
\min\!\Bigl\{1,\frac{\|g_k\|}{M\Delta_{\min}}\Bigr\}
\ge
\frac{\|g_k\|}{M\Delta_{\min}+ \|g_k\|}.
\tag{22}
\]
Inequality (22) implies
\[
\|g_k\|^{2}
\min\!\Bigl\{1,\frac{\|g_k\|}{M\Delta_{\min}}\Bigr\}
\;\ge\;
\frac{\|g_k\|^{3}}{M\Delta_{\min}+ \|g_k\|}.
\tag{23}
\]
Hence (21) gives
\[
\sum_{k\in\mathcal{S}} \frac{\|g_k\|^{3}}{M\Delta_{\min}+ \|g_k\|}
\;\le\;
\frac{2M}{\eta_1\kappa_{\mathrm{c}}}\bigl(f(x_0)-f_{\inf}\bigr).
\tag{24}
\]

\noindent\textbf{[Step 9]} (\emph{Bounding the number of iterations needed for
\(\|g_k\|\le\varepsilon\)}).  
If \(\|g_k\|>\varepsilon\) for all successful iterations, then each term in
(24) is at least
\[
\frac{\varepsilon^{3}}{M\Delta_{\min}+ \varepsilon}
\;:=\;c_{\varepsilon}>0.
\]
Consequently,
\[
|\mathcal{S}|\;c_{\varepsilon}
\;\le\;
\frac{2M}{\eta_1\kappa_{\mathrm{c}}}\bigl(f(x_0)-f_{\inf}\bigr),
\]
which yields the bound on the number of successful iterations
\[
|\mathcal{S}|
\;\le\;
\frac{2M}{\eta_1\kappa_{\mathrm{c}}c_{\varepsilon}}
\bigl(f(x_0)-f_{\inf}\bigr)
=
\mathcal{O}\!\Bigl(\frac{M^{2}(f(x_0)-f_{\inf})}
{\kappa_{\mathrm{c}}^{2}\eta_1^{2}\varepsilon^{2}}\Bigr).
\tag{25}
\]
Since each unsuccessful iteration contracts the radius by at most a factor
\(\gamma_{\mathrm{dec}}\) and the radius is lower‑bounded by \(\Delta_{\min}\),
the total number of unsuccessful steps is also bounded by a constant multiple
of \(|\mathcal{S}|\).  Thus the overall iteration count satisfies the same
\(\mathcal{O}(\varepsilon^{-2})\) bound.

\noindent\textbf{[Step 10]} (\emph{Limit point argument}).  
If the algorithm generated infinitely many iterates, (25) forces
\(\|g_k\|\to0\) along a subsequence of successful steps.  Because any
unsuccessful step leaves the point unchanged (\(x_{k+1}=x_k\)), the same limit
holds for the full sequence, establishing
\[
\liminf_{k\to\infty}\norm{\nabla f(x_k)}=0.
\qquad\square
\]

\section*{Rate Derivation (With Constants)}
Collecting the constants from the proof, the iteration complexity can be written
explicitly as
\[
N(\varepsilon)
\;\le\;
\Biggl\lceil
\frac{2M^{2}}{\eta_1\kappa_{\mathrm{c}}(M\Delta_{\min}+\varepsilon)\,\varepsilon^{2}}
\bigl(f(x_0)-f_{\inf}\bigr)
\Biggr\rceil .
\]
If we use the safe radius \(\Delta_{\min}=\Delta_{\mathrm{safe}}
=\frac{\kappa_{\mathrm{c}}M}{2L}\), the bound simplifies to
\[
N(\varepsilon)
\;\le\;
\Biggl\lceil
\frac{4L}{\eta_1\kappa_{\mathrm{c}}^{2}}
\frac{M}{\varepsilon^{2}}
\bigl(f(x_0)-f_{\inf}\bigr)
\Biggr\rceil .
\]

\section*{Special Cases}
\begin{itemize}[leftmargin=*]
\item \textbf{Exact subproblem solution.}
If \(p_k\) is the exact minimiser of (TR), then \(\kappa_{\mathrm{c}}=1\) and the
constants improve accordingly.
\item \textbf{Quadratic objective.}
When \(f\) itself is quadratic, the model is exact, \(\operatorname{ared}_k=
\operatorname{pred}_k\) and \(\rho_k=1\) for every iteration.  The algorithm
reduces to the classical trust‑region Newton method with a finite‑step
termination after at most \(n\) successful steps.
\item \textbf{Convex case.}
If \(f\) is convex, the same analysis holds, and the algorithm enjoys the
additional guarantee that every accepted step yields a monotone decrease of
\(f\).
\end{itemize}

\end{document}